\documentclass[11pt,letterpaper]{article}
\usepackage[T1]{fontenc}
\usepackage{lmodern,amsmath,amssymb,amsthm,booktabs}
\usepackage[margin=1in,top=.9in,bottom=.9in,headheight=14pt]{geometry}
\usepackage{microtype,xcolor,fancyhdr}
\usepackage[colorlinks=true,linkcolor=teal!65!black,citecolor=teal!65!black,urlcolor=teal!65!black]{hyperref}
\definecolor{ink}{RGB}{24,50,61}\definecolor{accent}{RGB}{29,104,103}
\pagestyle{fancy}\fancyhf{}\fancyhead[L]{\small\sffamily Prime series in arithmetic progressions}
\fancyhead[R]{\small\sffamily Residues and zero multiplicities}
\fancyfoot[C]{\small\thepage}\renewcommand{\headrulewidth}{.3pt}
\setlength{\parindent}{1.1em}\setlength{\emergencystretch}{1.3em}
\newtheorem{theorem}{Theorem}\newtheorem{corollary}[theorem]{Corollary}
\newcommand{\C}{\mathbb C}\newcommand{\R}{\mathbb R}\newcommand{\Z}{\mathbb Z}
\newcommand{\Res}{\operatorname{Res}}\newcommand{\supp}{\operatorname{supp}}
\newcommand{\ord}{\operatorname{ord}}\newcommand{\imag}{\operatorname{Im}}
\newcommand{\re}{\operatorname{Re}}\newcommand{\e}{\mathrm e}\newcommand{\dd}{\,\mathrm d}
\newcommand{\wt}{\widehat}\newcommand{\one}{\mathbf 1}
\hypersetup{pdftitle={Reading arithmetic weights from poles of a prime series},pdfauthor={OpenAI}}
\begin{document}
\thispagestyle{plain}
\noindent{\small\sffamily\color{accent} ANALYTIC NUMBER THEORY\hfill 26 NOVEMBER 2026}\par
\vspace{.55em}
\noindent{\LARGE\bfseries\color{ink}Reading arithmetic weights from poles\\[2pt] of a prime series}\par
\vspace{.4em}
\noindent{\large Power residues and local independence of Dirichlet zero multiplicities}\par
\vspace{.65em}
\noindent\textbf{Abstract.}
Assign a complex weight to each reduced residue class modulo a fixed integer. A short piece of a single vertical line of poles in the associated prime Dirichlet series determines precisely the weights on the classes that are powers of a given order. We prove an exact rank formula, uniformly over all orders and all sufficiently large heights. In the natural height coordinate of the zeros, the observation window has radius $(1/\pi+\varepsilon)\log\log\log |T|+O_{q,\varepsilon}(1)$. The proof handles cancellation at shared zeros by recovering every character coefficient from local multiplicities. We also classify the possible patterns of lines with poles and determine the minimum number of residue classes needed to realize any pattern.

\section{What one line of poles determines}
Fix $q\ge2$, let $G=(\Z/q\Z)^\times$, and give each $a\in G$ a weight $f(a)\in\C$. Write
\begin{equation}
 F_f(s)=\sum_{p\nmid q}\frac{f(p\bmod q)\log p}{p^s},\qquad
 G^m=\{a^m:a\in G\},\qquad U_m f(a)=f(a^m).
 \label{eq:series}
\end{equation}
The sum is over primes and converges absolutely for $\re s>1$. Our zero-location input is GRH for Dirichlet $L$-functions, including RH for $\zeta$. This is the scope of the linked Lean theorem \cite{lean}. All logarithms are natural.

\begin{theorem}[Local recovery]\label{thm:main}
Using Dirichlet GRH for the pole locations, the following assertions hold. The function $F_f$ continues meromorphically to $\re s>0$, with at most simple poles. Fix $b>1/\pi$. There are constants $C,T_0$ depending only on $q,b$ such that for every $m\ge1$ and $|T|\ge T_0$ the following holds. Observe the poles and residues on
\begin{equation}
  s=\frac{1/2+i\gamma}{m},\qquad
  |\gamma-T|\le H(T),\qquad H(T)=b\log\log\log|T|+C.
  \label{eq:window}
\end{equation}
For squarefree $m$, two weights give identical observations exactly when they agree on $G^m$. The linear observation map has rank $|G^m|$. For nonsquarefree $m$ it has rank zero. The constants are uniform in $f$ and $m$.
\end{theorem}

The observation is a finitely supported complex measure in the coordinate $\gamma$, with the residue as the mass at each pole. A removable singularity contributes zero. The radius in the $s$-plane is $H(T)/m$. At $m=1$ a single such segment determines every weight. At $m=2$ it determines exactly the weights on squares. The result allows arbitrary coincidences and multiplicities among the zeros.

\section{Separating the contributions of different powers}
Let $\widehat G$ denote the characters modulo $q$. Use the Fourier convention
\[
 c_\chi(g)=|G|^{-1}\sum_{a\in G}g(a)\overline{\chi(a)},\qquad
 g(a)=\sum_{\chi\in\widehat G} c_\chi(g)\chi(a).
\]
We write $L'/L(z,\chi)$ for the logarithmic derivative with respect to the first argument, evaluated at $z$. For $\re s>1$ the Euler product and M\"obius inversion give
\begin{align}
 F_f(s)
 &=-\sum_{k\ge1}\mu(k)\sum_{\chi\in\widehat G}
        c_\chi(f)\frac{L'}{L}(ks,\chi^k)\notag\\
 &=-\sum_{k\ge1}\mu(k)\sum_{\psi\in\widehat G}
        c_\psi(U_k f)\frac{L'}{L}(ks,\psi).
 \label{eq:inversion}
\end{align}
Indeed the coefficient of $p^{-js}$ for $p\nmid q$ is
$f(p^j)\log p\sum_{k\mid j}\mu(k)$, which leaves just $j=1$. The second equality follows by expanding $f(a^k)$ in characters. For any $\sigma_0>0$, the terms with $k\sigma_0\ge 2$ converge normally on $\re s\ge\sigma_0$. Bound each logarithmic derivative there by
$2^{-(k\sigma_0-2)}\sum_{n\ge2}\Lambda(n)/n^2$ and use
$\sum_\psi |c_\psi(U_kf)|\le\sqrt{|G|}\,\|f\|_\infty$.
The remaining finite sum is meromorphic with at most simple poles.

Let $\psi^*$ induce $\psi$ primitively, including the conductor-one character. At any point $\rho$ of the open critical strip put $n_\psi(\rho)=\ord_\rho L(\cdot,\psi^*)$, with value zero at a regular nonzero point. The omitted Euler factors have all their zeros on $\re z=0$. For $g:G\to\C$ put
\[
  Z_g(\rho)=\sum_\psi c_\psi(g)n_\psi(\rho).
\]
Under Dirichlet GRH, a nonreal pole of a $k$-term of \eqref{eq:inversion} in the right half-plane can occur only at $\re s=1/(2k)$. Thus for $\rho=1/2+i\gamma$ with $\gamma\ne0$ we have
\begin{equation}
 \Res_{s=\rho/m}F_f(s)=-\frac{\mu(m)}{m}Z_{U_m f}(\rho).
 \label{eq:res}
\end{equation}
The factor $1/m$ comes from changing the variable of a simple pole. On the real axis principal-character poles can meet central zeros from a different term, hence the restriction on $\gamma$. Formula \eqref{eq:res} gives an upper bound for the rank. The reverse bound requires ruling out cancellation throughout an observation window.

\section{Local independence of zero multiplicities}
We prove the analytic assertion in the entire open critical strip. In this form it is independent of GRH. Let $\mathcal Z_q$ be the set of distinct nontrivial zeros of the inducing primitive functions for $\widehat G$. For $\rho\in\mathcal Z_q$ write $u_\rho=(\rho-1/2)/i$, so that $\re u_\rho=\imag\rho$ and $|\imag u_\rho|<1/2$.

\begin{theorem}\label{thm:local}
Fix $q$ and $b>1/\pi$. In every window $|\imag\rho-T|\le b\log\log\log|T|+C_{q,b}$ with $|T|$ sufficiently large, the vectors
$\bigl(n_\psi(\rho)\bigr)_{\psi\in\widehat G}$ span $\C^{\widehat G}$. More precisely, there is a linear transformation $D_T$ of the weighted data $\{Z_g(\rho)\}$ in this window into $\C^G$ such that
\begin{equation}
 \|D_T Z_g-g\|_\infty\le\tfrac12\|g\|_\infty .
 \label{eq:norm}
\end{equation}
Any prescribed positive relative error in place of $1/2$ is possible by adjusting the transformation, $C_{q,b}$ and the height threshold.
\end{theorem}

\subsection*{Test functions that distinguish primes}
Choose a prime $p_a>2$ representing each $a\in G$, using Dirichlet's theorem. Only this finite list is used. Write $\lambda=\log p$ for $p=p_a$, fix $1/2<d<\pi b/2$, and define
\[
  k(u)=\exp[-\cosh(u/b)],\qquad
  K_R(u)=\frac1{2R}\int_{-R}^{R}k(u-v)\dd v,\qquad
  M=\int_{\R}k(u)\dd u>0.
\]
The averaging length $R>0$ controls frequency selectivity and is chosen independently of $T$. Use $\wt h(\xi)=\int h(u)\e^{-iu\xi}\dd u$. Shifting the contour by $-i d\,\operatorname{sgn}\xi$ gives
\begin{equation}
 |\wt k(\xi)|\le C_{b,d}\e^{-d|\xi|},\qquad
 \wt K_R(\xi)=\wt k(\xi)\frac{\sin(R\xi)}{R\xi}.
 \label{eq:ft}
\end{equation}
Indeed, $|k(x+iy)|=\exp[-\cos(y/b)\cosh(x/b)]$ and $\cos(d/b)>0$. Thus $\wt K_R(0)=M$ and $\wt K_R(\xi)\to0$ for each fixed $\xi\ne0$ as $R\to\infty$. For real $T$ put
\begin{align}
 h_{p,T}(u)&=\frac1{i\lambda}\frac{\mathrm d}{\mathrm du}
       \bigl[\e^{i\lambda u}K_R(u-T)\bigr],\notag\\
 \wt h_{p,T}(\xi)&=\frac{\xi}{\lambda}\,
        \e^{-i(\xi-\lambda)T}\wt K_R(\xi-\lambda).
 \label{eq:h}
\end{align}
In particular $\wt h_{p,T}(0)=0$ and $\wt h_{p,T}(\log p)=M$. Differentiation removes the main archimedean term. The phase in \eqref{eq:h} leaves the target prime unchanged as the window moves.
\subsection*{Recovering a character value}
For a character $\psi$ with primitive inducer $\eta$ of conductor $Q$ and parity $a_\eta\in\{0,1\}$, put
\[
 A_\eta(u)=\log(Q/\pi)+\re\frac{\Gamma'}{\Gamma}
                \left(\frac{1/2+a_\eta+iu}{2}\right).
\]
The Guinand--Weil formula in our Fourier convention reads \cite[(2.1), (2.6)]{gw}
\begin{align}
 \sum_\rho n_\psi(\rho) h(u_\rho)
   ={}&\one_{Q=1}\bigl(h(i/2)+h(-i/2)\bigr)
       +\frac1{2\pi}\int_{\R}h(u)A_\eta(u)\dd u\notag\\[-2pt]
   &-\frac1{2\pi}\sum_{n\ge2}\frac{\Lambda(n)}{\sqrt n}
       \left[\eta(n)\wt h(\log n)+\overline{\eta(n)}\wt h(-\log n)\right].
 \label{eq:gw}
\end{align}
The sum over $\rho$ runs over distinct points of $\mathcal Z_q$. The entire tests in \eqref{eq:h} decay in every closed horizontal strip of half-width less than $\pi b/2$, including one wider than $1/2$. The formula applies to the complex test by linearity.

Multiply \eqref{eq:gw} by $-2\pi\sqrt p/(M\log p)$. In the prime sum, the term $n=p$ at positive frequency gives exactly $\eta(p)=\psi(p)$. The absolute sum of all other contributions tends to zero as $R\to\infty$, uniformly in $T$ and $\eta$. Indeed, \eqref{eq:ft}--\eqref{eq:h} bound its summands, up to a constant depending on $p,b,d$, by
\[
 \frac{\Lambda(n)\log n}{\sqrt n}\,
       \exp\!\big(-d|\,{\pm}\log n-\log p|\big).
\]
Both majorants have convergent sums because $d>1/2$. At each remaining frequency the sinc factor tends to zero. Dominated convergence proves the uniform bound. This estimate includes powers of primes dividing $q$ when the primitive character has fewer missing Euler factors than $\psi$.

Fix $R$ after setting the desired arithmetic error. The two pole terms $h(\pm i/2)$ tend to zero as $|T|\to\infty$. For the gamma term, integration by parts gives
\[
 \int_{\R} h_{p,T}(u)A_\eta(u)\dd u
  =-\frac1{i\lambda}\int_{\R} \e^{i\lambda u} K_R(u-T) A_\eta'(u)\dd u=o(1).
\]
Stirling's formula gives $A_\eta'(u)=O((1+|u|)^{-1})$ at infinity, and the derivative is bounded on the real axis. Translate by $T$ and apply dominated convergence to $K_R$. Both signs of $T$ are covered.

\subsection*{Discarding zeros outside the window}
On $|\imag u|\le1/2$, for constants $c,C'>0$ depending on $p,b,R$, we have
\begin{equation}
 |h_{p,T}(u)|\le C'\exp\!\left[-c\exp\bigl(|\re u-T|/b\bigr)\right].
 \label{eq:decay}
\end{equation}
Use the strip bound for $k$ and the bounded translation $|v|\le R$. Differentiation introduces only a hyperbolic sine, absorbed by decreasing $c$. The classical local zero count with multiplicities is $O_q(\log(3+|t|))$ on $t\le\imag\rho\le t+1$, by the Riemann--von Mangoldt formula. Grouping both tails into unit intervals starting at distance $H$ gives, for $H\ge1$,
\begin{equation}
 \sum_{|\imag\rho-T|>H}n_\psi(\rho)|h_{p,T}(u_\rho)|
 \ll_{q,b,R,p}(1+\log(3+|T|))\exp[-c'\exp(H/b)].
 \label{eq:tail}
\end{equation}
At distance $H+j$ the logarithmic count is $O_q(\log(3+|T|)+\log(3+H+j))$.
Keep half of $c\exp((H+j)/b)$ to produce the factor $\exp[-c'\exp(H/b)]$ and use the other half to sum the logarithmic factor uniformly in $H$. For $H(T)$ in \eqref{eq:window}, $\exp(H/b)=\e^{C_{q,b}/b}\log\log |T|$. Increasing $C_{q,b}$ makes \eqref{eq:tail} tend to zero. This tail estimate explains the third iterated logarithm.

\begin{proof}[Completion of Theorem \ref{thm:local}]
Define the $a$ coordinate of the transformation by
\[
 (D_T Z_g)(a)=\frac{-2\pi\sqrt{p_a}}{M\log p_a}
     \sum_{|\imag\rho-T|\le H(T)}h_{p_a,T}(u_\rho) Z_g(\rho).
\]
For a single character $g=\psi$, the preceding estimates make the error from $\psi(a)$ less than any chosen $\delta>0$, simultaneously for all $a,\psi$. First choose $R$ for the frequency sum, then $C_{q,b}$ for the tail, then the height threshold. For general $g$, Parseval and Cauchy--Schwarz give
$\sum_\psi|c_\psi(g)|\le\sqrt{|G|}\|g\|_\infty$. Taking $\delta<1/(2\sqrt{|G|})$ proves \eqref{eq:norm}. If all the observed multiplicity combinations vanish, this inequality forces $g=0$. Fourier inversion gives full column rank, or equivalently the asserted span.
\end{proof}

\section{Exact inversion and the boundary of continuation}
\begin{proof}[Proof of Theorem \ref{thm:main}]
For squarefree $m$, multiply the observed residues by $-m/\mu(m)$. By \eqref{eq:res} these values are exactly $Z_{U_m f}(\rho)$ on the window. Increase $T_0$ so that the window misses the real axis. By Theorem \ref{thm:local} the data uniquely determine $U_m f$. In particular $f(y)$ is known for every $y\in G^m$, by choosing an $m$th root in $G$. Conversely all the residues factor through this restriction. Every function on $G^m$ has an extension to $G$, so the image has dimension $|G^m|$. For nonsquarefree $m$ the residue is zero by \eqref{eq:res}.
\end{proof}

For exact recovery of complex weights one can calibrate a finite matrix. Its $a'$ column is $D_T Z_{\one_{\{a'\}}}$, computed from the known family of Dirichlet functions for $q$. By \eqref{eq:norm} this matrix lies within $1/2$ of the identity in the induced supremum norm. Its inverse has norm at most $2$. Applying the inverse to $D_T Z_{U_m f}$ recovers $U_m f$ exactly. The bound concerns errors in normalized probe outputs. Computing certified zero lists and propagating errors in the zero locations require further numerical estimates. For integer weights bounded by a known $B\ge1$, one may choose relative error less than $1/(2B)$ and round the real part of each probe output directly.

\begin{corollary}[Lines with poles and the natural boundary]\label{cor:boundary}
The line $\re s=1/(2m)$ contains a nonreal pole of $F_f$ exactly when $m$ is squarefree and $\supp(f)\cap G^m\ne\varnothing$. Each such line has poles in every sufficiently high window \eqref{eq:window} in either direction. If $f\ne0$, the imaginary axis is a natural boundary for meromorphic continuation.
\end{corollary}
\begin{proof}
Compare $f$ with zero in Theorem \ref{thm:main}. If $E$ is the exponent of $G$, the power map is an automorphism for every prime $m\nmid E$. For $f\ne0$ each such $m$ is an index with poles. For real $t\ne0$ take $T=mt$. One obtains poles $s_m$ with
\[
  \re s_m=1/(2m),\qquad
  |\imag s_m-t|\le \bigl(b\log\log\log|mt|+C_{q,b}\bigr)/m.
\]
For $t=0$ use a single sufficiently large fixed $T$ as $m$ varies. In each case distinct poles accumulate at $it$, which precludes a meromorphic extension across that point.
\end{proof}

\section{Which patterns of poles can occur?}
Let $P$ be the set of prime divisors of the exponent of $G$. Put $S(m)=\{\ell\in P:\ell\mid m\}$ and
$B(y)=\{\ell\in P:y\in G^\ell\}$. A family $\mathcal D\subseteq 2^P$ is downward closed if $S\in\mathcal D$ and $S'\subseteq S$ imply $S'\in\mathcal D$.

\begin{theorem}[Prescribing the lines]\label{thm:design}
On squarefree indices, the possible patterns of lines with nonreal poles are precisely the downward closed families of sets $S(m)$. The smallest size of $\supp(f)$ realizing a given family equals its number of inclusion-maximal members. The empty family has minimum support zero.
\end{theorem}
\begin{proof}
Decompose $G$ into its Sylow subgroups. On the $\ell$-part, an exponent prime to $\ell$ acts invertibly. Thus for squarefree $m$,
$G^m=\bigcap_{\ell\in S(m)}G^\ell$. A supported point $y$ makes precisely the subsets of $B(y)$ visible. Their union is downward closed. Conversely every $B\subseteq P$ is $B(y)$ for some $y$. On the Sylow $\ell$-subgroup choose the identity for $\ell\in B$, and choose an element outside its $\ell$th powers for $\ell\notin B$. Such an element exists in every nontrivial finite abelian $\ell$-group. Choose one such $y$ for each maximal member of the desired family and assign it weight one. Each maximal member must equal $B(y)$ for some supported $y$, so each requires its own point.
\end{proof}

For a family $\mathcal D$ the relative natural density among squarefree indices is
\begin{equation}
  \sum_{S\in\mathcal D}
     \prod_{\ell\in S}\frac1{\ell+1}
     \prod_{\ell\in P\setminus S}\frac{\ell}{\ell+1}.
  \label{eq:density}
\end{equation}
Among the $\ell^2-1$ allowable residues modulo $\ell^2$, precisely $\ell-1$ are divisible by $\ell$. These finite conditions are independent by the Chinese remainder theorem. Squarefreeness at the remaining primes multiplies each density by the same Euler product. To see this, truncate at $L$ and bound the upper density of the discarded multiples of prime squares by $\sum_{p>L}p^{-2}$. Multiplying \eqref{eq:density} by $6/\pi^2$ gives the density among all positive integers.

\subsection*{Two series modulo 31}
The group has order $30$ and generator $3$, so $3^j\in G^m$ precisely when $\gcd(m,30)\mid j$. Write $\delta_a$ for the indicator of the residue $a$. In the table $m$ is squarefree and $d=\gcd(m,30)$. The last two columns record whether a line contains nonreal poles.
\begin{center}
\begin{tabular}{r r c c}
\toprule
$d$ & rank & $f=\delta_3$ & $f=\delta_9-\delta_{27}$ \\
\midrule
1 &30&yes&yes\\
2 &15&no&yes\\
3 &10&no&yes\\
5 &6&no&no\\
6 &5&no&no\\
10&3&no&no\\
15&2&no&no\\
30&1&no&no\\
\bottomrule
\end{tabular}
\end{center}
For $\delta_3$ the relative density is $(2/3)(3/4)(5/6)=5/12$. The signed weight gives $+1$ at $9=3^2$ and $-1$ at $27=3^3$. Among squarefree indices its lines carry poles precisely when $5\nmid m$ and $6\nmid m$. The relative density is $(5/6)(1-1/(3\cdot4))=55/72$. This pattern needs two supported classes by Theorem \ref{thm:design}, since its maximal sets are $\{2\}$ and $\{3\}$. In particular, poles occur on the lines $\re s=1/4$ and $1/6$, and the series is holomorphic at every nonreal point of $\re s=1/12$.

\subsection*{Reading two residues modulo four}
Modulo four the character family consists of the principal character, induced by $\zeta$, and the odd character with Dirichlet beta function $\beta$. At $m=1$, suppose one observes a simple beta zero unshared by zeta, and a simple zeta zero unshared by beta. After changing the sign of the two residues as in \eqref{eq:res}, the measurements are
\[
 z_\beta=\tfrac12(f(1)-f(3)),\qquad z_\zeta=\tfrac12(f(1)+f(3)).
\]
Consequently $f(1)=z_\zeta+z_\beta$ and $f(3)=z_\zeta-z_\beta$. For the signed weight $(1,-1)$ the zeta zero disappears from $F_f$, while the beta zero gives residue $-1$. At $m=2$, both unit classes square to $1$, and only $f(1)$ is recoverable.

The accompanying numerical demonstration finds the indicated beta and zeta roots at heights $12.9880980123124\ldots$ and $14.1347251417347\ldots$, respectively. It checks residuals below $10^{-32}$ and evaluates the other function and the Hardy derivative at each root. These floating-point calculations illustrate the two rows above. Theorem~\ref{thm:main} uses exact multiplicities and allows shared zeros.

\section{A consequence for two individual Dirichlet functions}
There is a distinction between detecting that two multiplicities differ and specifying which one is larger. For characters of the same primitive conductor, positivity gives both directions in a single window.

\begin{corollary}\label{cor:directed}
Assume Dirichlet GRH. Let $\chi_1,\chi_2$ be distinct primitive Dirichlet characters of the same conductor. For every $b>1/\pi$ and all sufficiently large $|T|$, a window of radius $b\log\log\log|T|+O_{\chi_1,\chi_2,b}(1)$ contains a zero where $n_{\chi_1}>n_{\chi_2}$ and another where $n_{\chi_2}>n_{\chi_1}$.
\end{corollary}
\begin{proof}
Put $v(\rho)=n_{\chi_1}(\rho)-n_{\chi_2}(\rho)$ and choose a prime $p$ at which the characters differ. In the difference of their explicit formulas use the two undifferentiated tests $K_R(u-T)$ and $\e^{iu\log p}K_R(u-T)$. The conductors agree, so their terms $\log(Q/\pi)$ cancel exactly. The difference of the gamma contributions to $A_\eta$ is bounded and tends to zero at either infinity by Stirling's formula. Its integral against either test is $o_{|T|\to\infty}(1)$ for fixed $R$. Neither primitive function has a pole, since conductor one has only one primitive character.

At frequency zero there is no prime-power term. At positive frequency $\log p$ the normalized Fourier peak is $M$. The same dominated convergence as in \eqref{eq:ft}--\eqref{eq:gw}, now without the factor $\xi/\lambda$, and the tail bound \eqref{eq:tail} apply to both tests. Write
\[
 d_p=\frac{M\log p}{2\pi\sqrt p}|\chi_1(p)-\chi_2(p)|>0 .
\]
Choose $R$ large, then the window constant and starting height large, so that in the window $I$ the untwisted sum has modulus $< d_p/3$ and the twisted sum has modulus $>2d_p/3$. Here the sums are $\sum_{\gamma\in I} v(1/2+i\gamma)K_R(\gamma-T)$ and the same sum with the factor $\e^{i\gamma\log p}$. We have used Dirichlet GRH to make the ordinates real in the tests. Since $K_R$ is positive on $\R$ and the twisting factor has modulus one, a single sign of $v$ throughout $I$ would bound the second modulus by the first. Both signs must occur.
\end{proof}

\section{Sources, scope and further questions}
\paragraph{The formal input.}
The readable edition of \cite{lean} exposes five statements in \texttt{grh/src/Main.lean}. They give Mathlib RH, its expanded zeta formulation, Dirichlet nonvanishing to the right of the critical line with the principal pole excepted, the critical-strip theorem for all Dirichlet characters, and its primitive specialization. We use Dirichlet zero location, including the zeta case. The reading guide explains how short forms, a spanning estimate and a lower bound for integral forms yield a one-sided bound for a smoothed sum of $\Lambda(n)(1-\chi_d(n))$ in a progression. An explicit-formula argument carries that bound to a zero-free half-plane. This reduction is part of the iterated theta-lift approach of Alp\"oge and Claude \cite{alpoge}. Here zero location separates the prime powers, and local zero data identifies the character coefficients. The RH article \cite{wiki} gives the critical-line and prime-distribution context. Rubinstein reports an independent replay of the compact rewrite \cite{replay}. We inspected the source statements and used the published formal theorem as input.

\paragraph{The ProveIt comparison.}
The index at the pinned revision \cite{proveit} describes 61 packets archived without repository-level verification. A close comparison is \emph{Recovering both arithmetic parameters from a finite signed zero list}, in \texttt{grh-torsion-and-zero-data/recovery.tex}. It treats a particular signed combination of untwisted and quadratic-twisted multiplicities, with a height bound $O(d^2\log^3(2d))$ for recovering unknown parameters $q,d$. The packet \texttt{exponent\_seven\_and\_grh\_transfers} also treats a two-sided dyadic band. In the present problem $q$ is known, the weight is arbitrary and complex, and the output is its exact restriction to $G^m$ from any sufficiently high, short segment of the $m$th line.

\paragraph{Earlier results.}
Scalar prime-zeta continuation is classical \cite{froberg}. Mathar gives character-prime M\"obius inversion \cite{mathar}. Alberts develops the general factorization method for Frobenian Euler products \cite{alberts}. Kaczorowski and Perelli prove global independence of multiplicity vectors in a general setting \cite[Lemma 2]{kp}. Banks distinguishes pairs of Dirichlet zero multisets in logarithmic windows with explicit conductor dependence \cite{banks}. Srinivas obtains a prescribed direction of multiplicity excess in double-logarithmic windows for Selberg-class pairs under a degree ordering \cite{srinivas}. Corollary~\ref{cor:directed} gives the shorter scale for the fixed equal-conductor Dirichlet case. Double-exponential kernels have a history in short-interval problems \cite{ivic}.

Our proposed contribution is simultaneous full rank in every triple-logarithmic window for a fixed character family, its exact power-image rank for periodic prime series, and the realization theorem. A targeted comparison found no earlier statement of this combination. Historical priority remains open to further bibliographic checking. The new proofs are analytic. The accompanying code checks finite arithmetic identities and illustrates local rank numerically modulo four.

\paragraph{Further questions.}
Does a fixed-radius window suffice at a fixed modulus? How small can the additive constant and starting height be made as explicit functions of $q$? How much error can be allowed in the locations and residues when using certified interval data? For a nonabelian class function, what replaces the power-image rank when Artin factors share constituents or have poles? Optimality of the triple-logarithmic scale and of the coefficient $1/\pi$ is open in this analysis.

\clearpage
\begin{thebibliography}{99}\small
\bibitem{alpoge} L. Alp\"oge, \href{https://alpo.ge/remark.pdf}{Announcement of the Lean proof of RH and Dirichlet GRH}, November 2026. Author's announcement and outline.
\bibitem{lean} L. Sernau and collaborators, \emph{GRH for Dirichlet $L$-functions, in Lean 4}, revision \texttt{72fead1}. \href{https://github.com/lsernau/grh-lean/tree/72fead102f663d5c7a407964a13dc0cf9d977207}{Pinned repository}, \href{https://github.com/lsernau/grh-lean/blob/72fead102f663d5c7a407964a13dc0cf9d977207/grh/src/Main.lean}{theorem declarations} and \href{https://github.com/lsernau/grh-lean/blob/72fead102f663d5c7a407964a13dc0cf9d977207/grh/GUIDE.md}{reading guide}.
\bibitem{proveit} V. Reshetnikov, \emph{ProveIt}, GRH research collection, revision \texttt{61e72d7}. \href{https://github.com/VladimirReshetnikov/ProveIt/tree/61e72d75affa1084cf3ac4f66d199f54efde05fe/Analysis/GeneralizedRiemannHypothesis}{Index} and \href{https://github.com/VladimirReshetnikov/ProveIt/blob/61e72d75affa1084cf3ac4f66d199f54efde05fe/Analysis/GeneralizedRiemannHypothesis/Research/grh-torsion-and-zero-data/recovery.tex}{finite signed-zero recovery}.
\bibitem{wiki} Wikipedia contributors, \href{https://en.wikipedia.org/wiki/Riemann_hypothesis}{\emph{Riemann hypothesis}}, consulted 26 November 2026.
\bibitem{replay} M. Rubinstein, \href{https://www.math.uwaterloo.ca/~mrubinst/alpoge_claude_grh/}{\emph{Alp\"oge/Claude proof of the Generalized Riemann Hypothesis}}, verification reports, November 2026.
\bibitem{gw} A. Chirre, V. J. Pereira J\'unior and D. de Laat, \emph{Primes in arithmetic progressions and semidefinite programming}, \href{https://arxiv.org/abs/2005.02393}{arXiv:2005.02393}, equations (2.1) and (2.6).
\bibitem{froberg} C.-E. Fr\"oberg, \emph{On the prime zeta function}, BIT 8 (1968), 187--202. \href{https://doi.org/10.1007/BF01933420}{doi:10.1007/BF01933420}.
\bibitem{mathar} R. J. Mathar, \emph{Table of Dirichlet $L$-series and prime zeta modulo functions for small moduli}, \href{https://arxiv.org/abs/1008.2547}{arXiv:1008.2547}, Section 3.
\bibitem{alberts} B. Alberts, \emph{Explicit analytic continuation of Euler products}, Essential Number Theory 5 (2026), 49--112. \href{https://doi.org/10.2140/ent.2026.5.49}{doi:10.2140/ent.2026.5.49}.
\bibitem{kp} J. Kaczorowski and A. Perelli, \emph{Functional independence of the singularities of a class of Dirichlet series}, Amer. J. Math. 120 (1998), 289--303. \href{https://doi.org/10.1353/ajm.1998.0012}{doi:10.1353/ajm.1998.0012}.
\bibitem{banks} W. D. Banks, \emph{Comparing zeros of distinct Dirichlet $L$-functions}, \href{https://arxiv.org/abs/2302.07073}{arXiv:2302.07073}, Theorems 1.1--1.2.
\bibitem{srinivas} K. Srinivas, \emph{Distinct zeros of functions in the Selberg class}, Acta Arith. 103 (2002), 201--207. \href{https://doi.org/10.4064/aa103-3-1}{doi:10.4064/aa103-3-1}.
\bibitem{ivic} A. Ivi\'c, \emph{On mean value results for the Riemann zeta-function in short intervals}, Hardy--Ramanujan J. 32 (2009). \href{https://hrj.episciences.org/164/pdf}{Full text}.
\end{thebibliography}
\end{document}
